\section{Study Design}

\subsection{Stage 1: matched-context feasibility}

The first study uses a randomized within-participant crossover of S, E, and A. Condition order is counterbalanced or randomized, and tasks are matched on nominal cognitive rule and challenge level. The purpose is to establish task equivalence, tolerability, movement exposure, telemetry quality, and the magnitude of context effects.

This stage is not a test of clinical efficacy.

\subsection{Stage 2: mechanism experiment}

After acceptable matching is demonstrated, repeated eligible blocks can be randomized among S, E, and A. The primary context-specific contrasts are \(E-S\) and \(A-E\). Condition order, fatigue, session number, challenge type, and prior exposure are modeled or balanced according to a preregistered design.

\subsection{Stage 3: contextual-transfer experiment}

Only after the proximal contrasts are understood should a longitudinal study test whether training across varied contexts improves independent transfer. One possible design compares a fixed-context training schedule with a context-varied schedule while holding total training dose, cognitive content, reward structure, and challenge exposure constant.

\subsection{BYOD hardware qualification}

The embodied/AR programme uses family-owned Android phones under a stricter module-specific qualification layer than the core intervention. Qualification must establish required camera/AR capability, acceptable tracking/runtime behavior, relevant inertial-sensor availability, display performance, and any haptic requirements before condition assignment. Device capability class and technical failures are recorded so that context effects are not silently confounded with hardware failure. A phone that qualifies for the core intervention may still be ineligible for the AR companion module.

\subsection{AR implementation fidelity}

The research build logs AR tracking state, spatial-anchor failures, frame-rate degradation, device temperature events, and safety interruptions. AR blocks with technical failure are classified using predefined rules rather than removed because performance was poor.
